% Puissance dissipée par effet Joule pendant la charge, pour deux résistances (corrigé du défi D26)
\begin{tikzpicture}
\begin{axis}[otMesures,otAxesBout,otLecture,width=104mm,height=54mm,xmin=0,xmax=0.065,ymin=0,ymax=11,
  xtick={0,0.01,...,0.06},ytick={0,2,...,10},
  xticklabel style={/pgf/number format/fixed,/pgf/number format/precision=2},
  scaled x ticks=false,
  xlabel={$t$ (s)},ylabel={$P_\mathrm{J}$ (W)}]
\addplot[otNavy,line width=1pt,fill=otNavy,fill opacity=0.12,domain=0:0.06,samples=120] {10*exp(-2*x/0.01)} \closedcycle;
\addplot[otRed!80!black,line width=1pt,fill=otRed,fill opacity=0.10,domain=0:0.06,samples=120] {5*exp(-2*x/0.02)} \closedcycle;
\node[pcLegende,anchor=west,text=otNavy] at (axis cs:0.004,8.2) {$R=\qty{10}{\ohm}$ : $\tau=\qty{10}{ms}$};
\node[pcLegende,anchor=west,text=otRed!80!black] at (axis cs:0.017,3.6) {$R=\qty{20}{\ohm}$ : $\tau=\qty{20}{ms}$};
\node[pcLegende,anchor=west] at (axis cs:0.032,6.5) {même aire sous les deux courbes :};
\node[pcLegende,anchor=west] at (axis cs:0.032,5.5) {$W_\mathrm{J}=\frac12\,C\,E^2=\qty{0,05}{J}$};
\end{axis}
\end{tikzpicture}
\par\smallskip\legende{Avec une résistance double, la puissance dissipée est deux fois plus faible, mais pendant deux fois plus longtemps : l'énergie perdue ne dépend pas de $R$.}
